% purpose: 说明储电量轨迹被上下界夹成一条走廊，单段涨落幅度受功率上限限制且充放不对称，
%          日内首尾储电量还必须相等
% topology: geometry   reading_direction: coordinate-based
% node_roles: 上下界、可行走廊、可达斜率锥、首尾闭合线、示例轨迹
% visual_encoding: 虚线为界、粗实线为轨迹、点线为闭合参考，斜率由锥的两条边表示
% style_family: restrained Nature-style line art   print_mode: colorblind-safe   final_width: text
\documentclass[border=3pt,11pt]{standalone}
\usepackage{ctex}
\usepackage{amsmath}
\usepackage{tikz}
\usetikzlibrary{arrows.meta}
\definecolor{mhblue}{HTML}{0072B2}
\begin{document}
\begin{tikzpicture}[
  bound/.style={draw=mhblue,thin,dash pattern=on 3pt off 2pt},
  traj/.style={draw=mhblue,line width=1.1pt},
  tie/.style={draw=mhblue,thin,dash pattern=on 0.7pt off 1.6pt},
  cone/.style={draw=mhblue,line width=0.6pt,dash pattern=on 2pt off 1.5pt},
  axis/.style={draw=mhblue,line width=0.7pt,-{Stealth[length=4pt]}},
  lead/.style={draw=mhblue,line width=0.5pt},
  lb/.style={font=\small,inner sep=1.2pt}]

% ---- 坐标轴与时刻刻度 ----
\draw[axis] (0,0) -- (9.4,0);
\draw[axis] (0,0) -- (0,5.7);
\node[lb] at (9.3,0.285) {\(t\)};
\node[lb] at (-0.5,5.62) {\(E_t\)};
\foreach \x/\s in {0/0:00,2.25/6:00,4.5/12:00,6.75/18:00,9/24:00}{
  \draw[draw=mhblue,line width=0.6pt] (\x,0) -- (\x,-0.12);
  \node[lb] at (\x,-0.48) {\s};}

% ---- 储电量上下界：中间即可行走廊 ----
\draw[bound] (0,0.9) -- (9.2,0.9);
\draw[bound] (0,4.5) -- (9.2,4.5);
\node[lb,anchor=west] at (9.28,0.9) {\(E_{\min}\)};
\node[lb,anchor=west] at (9.28,4.5) {\(E_{\max}\)};

% ---- 示例轨迹：谷时段充、峰时段放，首尾同高 ----
\draw[traj] (0,2.2) -- (1.0,3.0) -- (1.8,3.6) -- (2.6,3.1) -- (3.4,2.3)
  -- (4.2,1.6) -- (5.0,2.4) -- (5.8,3.3) -- (6.6,4.0) -- (7.4,3.2)
  -- (8.2,2.5) -- (9.0,2.2);
\fill[black] (0,2.2) circle (2.4pt);
\fill[black] (9,2.2) circle (2.4pt);

% ---- 首尾闭合参考线 ----
\draw[tie] (0,2.2) -- (9.2,2.2);
\node[lb,fill=white,anchor=west] at (7.15,2.05) {\(E_0=E_{144}\)};

% ---- 单段可达斜率锥：充电与放电的最大涨落不对称 ----
\draw[cone] (0.55,1.55) -- (1.75,2.55);
\draw[cone] (0.55,1.55) -- (1.75,0.32);
\draw[cone] (1.75,2.55) -- (1.75,0.32);
\fill[black] (0.55,1.55) circle (1.9pt);
\node[lb,fill=white,anchor=west] at (1.85,2.62) {单段最大充电涨幅 \(\eta_c\bar{c}\)};
\node[lb,fill=white,anchor=west] at (1.85,0.35) {单段最大放电跌幅 \(\bar{s}/\eta_d\)};

% ---- 走廊标注 ----
\node[lb,fill=white,anchor=west] at (3.1,5.15) {可行走廊：轨迹须逐段落在上下界之间};
\draw[lead] (3.05,5.05) -- (2.35,4.55);

\end{tikzpicture}
\end{document}
